\clearpage\SourceLeaf{450}%
\begin{center}\begingroup\fontsize{8}{9}\selectfont
\setlength{\tabcolsep}{4pt}
\begin{tabular}{r|r|p{170pt}|r}\hline
\(+\) südlich&\(+\) westlich&&Nr.\\\hline
\(-242773.3\)&\(+36827.3\)&Drochtersen&\((1)\)\\
\(-235431.6\)&\(+29726.7\)&Colmar&\((1)\)\\
\(-245551.2\)&\(+90763.0\)&Cappeln&\((1)\)\\
\(-245775.5\)&\(+87628.9\)&Midlum&\((1)\)\\
\(-245821.6\)&\(+24600.1\)&Neuendorf&\((1)\)\\
\(-246532.7\)&\(+67702.1\)&Oster Ilienworth&\((1)\)\\
\(-246892.5\)&\(+76614.2\)&Wanna&\((1)\)\\
\(-247388.5\)&\(+63676.7\)&Bülkau&\((1)\)\\
\(-247703.2\)&\(+36752.9\)&Krautsand&\((1)\)\\
\(-250722.6\)&\(+69097.6\)&Neuenkirchen&\((1)\)\\
\(-250784.2\)&\(+65963.9\)&Osterbruch&\((1)\)\\
\(-251343.0\)&\(+34554.1\)&Glückstadt kleiner Thurm (Dänische Station)&\((1)\)\\
\(-251399.8\)&\(+30139.2\)&Herzhorn&\((1)\)\\
\(-251523.1\)&\(+11052.9\)&Barmstedt&\((1)\)\\
\(-251573.4\)&\(+34126.7\)&Glückstadt Kirchthurm&\((1)\)\\
\(-252146.1\)&\(+40748.4\)&Hammelvörden&\((1)\)\\
\(-253330.1\)&\(+78043.0\)&Lüdingworth&\((1)\)\\
\(-254110.4\)&\(+21378.6\)&Horst&\((1)\)\\
\(-254115.7\)&\(+68778.6\)&Otterndorf&\((1)\)\\
\(-255131.8\)&\(+28321.4\)&Süderau&\((1)\)\\
\(-255141.2\)&\(+62459.0\)&Belum&\((1)\)\\
\(-255213.6\)&\(+83123.0\)&Franzenburg&\((1)\)\\
\(-255856.6\)&\(+43050.3\)&Freiburg&\((1)\)\\
\hline\end{tabular}\endgroup\end{center}
\clearpage
\begin{center}\begingroup\fontsize{8}{9}\selectfont
\setlength{\tabcolsep}{4pt}
\begin{tabular}{r|r|p{170pt}|r}\hline
\(+\) südlich&\(+\) westlich&&Nr.\\\hline
\(-255908.9\)&\(+33800.4\)&Borsfleth&\((1)\)\\
\(-256062.6\)&\(+84523.8\)&HartmannsPlatz 1 bei Altenwalde&\((1)\)\\
\(-256062.7\)&\(+84104.8\)&Altenwalde&\((1)\)\\
\(-256069.1\)&\(+77120.7\)&Altenbruch Spitze 1.&\((1)\)\\
\(-256075.2\)&\(+77119.1\)&Altenbruch Spitze 2.&\((1)\)\\
\(-256894.2\)&\(+29816.6\)&Crempe&\((1)\)\\
\(-257722.7\)&\(+21032.0\)&Hohenfelde&\((1)\)\\
\(-257778.1\)&\(+35721.5\)&Wevelsfleth&\((1)\)\\
\(-257996.5\)&\(+80031.2\)&Groden&\((1)\)\\
\(-258998.8\)&\(+15610.9\)&Hörnerkirchen&\((1)\)\\
\(-259523.9\)&\(+26920.1\)&Neuenbrook&\((1)\)\\
\(-259779.2\)&\(+40370.0\)&Brockdorf&\((1)\)\\
\(-259854.2\)&\(+81878.1\)&Ritzebüttel Giebelstange 1&\((1)\)\\
\(-260086.9\)&\(+33836.4\)&Neuenkirchen&\((1)\)\\
\(-261494.4\)&\(+81196.7\)&Cuxhaven Leuchtthurm&\((1)\)\\
\(-261725.0\)&\(+34726.2\)&Beienfleth&\((1)\)\\
\(-262754.7\)&\(+83558.8\)&Döse&\((1)\)\\
\(-263188.1\)&\(+45161.9\)&St. Margareth&\((1)\)\\
\(-266552.1\)&\(+37341\,8\)&Wilster&\((1)\)\\
\(-266569.396\)&\(+95074.242\)&Neuwerk Leuchtthurm Cent.&\((1)\)\\
\(-269380.6\)&\(+14702.3\)&Kellinghusen&\((1)\)\\
\(-270223.4\)&\(+61150.8\)&Marne&\((1)\)\\
\hline\end{tabular}\endgroup\end{center}

Zur Erläuterung der Bedeutung der Coordinaten ist folgendes zu bemerken.

Will man sich nur im Allgemeinen einen Begriff davon machen, so kann man dieselben so ansehen, dass die erste Zahl anzeigt, wie viel der betreffende Ort südlich (beim \(+\) Zeichen), oder nördlich (beim \(-\) Zeichen) von der Göttinger Sternwarte liegt, die zweite Zahl hingegen, wie viel westlich (bei \(+\)) oder östlich (bei \(-\)).

Es ist aber dabei schon die Krummung der Erdoberfläche dergestalt berücksichtigt, dass bei Auftragung dieser Coordinaten auf eine ebene Fläche das Bild ein \textit{conformes}, d. i. in den kleinsten Theilen ähnliches wird. Das Nähere darüber enthalten meine geodätischen Abhandlungen zum Theil schon jetzt, und spätere Abhandlungen werden dies noch ausführlicher entwickeln.

Der genaue Anfangspunkt der Coordinaten in der Sternwarte ist übrigens der Mittelpunkt der Achse des Reichenbachschen Meridiankreises.

Als Einheit der Coordinaten ist diejenige Lineargrösse gewählt, die nach der besten im Jahr 1821 vorhandenen Kenntniss als der zehnmillionste Theil des Quadranten des Erdmeridians gelten konnte, nemlich die Länge von 443,307885 pariser Linien, was etwas, obwohl nur sehr wenig, von dem sogenannten legalen französischen Meter verschieden ist. Dies letztere war nemlich bekanntlich festgesetzt zu 443,296 pariser Linien. Obgleich in späterer Zeit (seit 1821) noch neuere Bestimmungen des zehnmillionten Theils des Erdmeridianquadranten gewonnen sind und zwar immer entschieden grösser als das eben angeführte gesetzliche Meter), so habe ich doch vorgezogen, bei der einmal von mir gewählten Einheit zu bleiben, da man jede einzelne Zahl leicht in jede beliebige andre Einheit umsetzen kann, zu welcher das Verhältniss einmal bekannt ist.
